% Abstract: at most 150 words (wc -w on the text without comments). Numbers follow paper/figures/NUMBERS.md.
\begin{abstract}
Narrowing the label set before prompting a large language model (LLM) is a cheap aid to in-context text classification.
Conformal In-Context Learning (CICLe) does this with a small classifier calibrated by class-conditional conformal prediction.
Earlier evidence for it came from a protocol with unlisted labels, unequal example pools and unseeded conformal sets, under which up to 65\% of outputs did not parse.
We correct it and evaluate CICLe on five English datasets with eight open LLMs and paired bootstrap tests.
Narrowing gives a small, reliable gain over few-shot prompting, 0.5 to 1.8 macro-F1 points, and 11 to 45\% shorter prompts.
How one narrows matters.
When the labelled pool is long-tailed, class-conditional sets keep their coverage while top-$k$, probability-mass and marginal sets lose 3 to 13 points.
The gain is largest when label names mean nothing.
With 1,600 balanced labels a fine-tuned encoder beats the 3B to 8B pipelines, which win when the pool is skewed or small.
\end{abstract}
